\documentclass{article}
\usepackage{amsmath,amssymb}
\usepackage[vocab={integralequations}]{leanify}
\newtheorem{lemma}{Lemma}
\begin{document}

\section{Notation}

For all $x,t\in\mathbb{R}$, let $K(x,t)\coloneqq1$.
For every real number $x$, let $e(x)\coloneqq\exp(x)$.
For every real number $x$, let $o(x)\coloneqq1$.
For every real number $x$, let $c(x)\coloneqq2$.
For every real number $x$, let $z(x)\coloneqq0$.
For every real number $x$, let $r(x)\coloneqq2\sqrt{x}$.
We write $K_2$ for the second iterated kernel of $K$ on $[0,1]$.
We write $T(y,x)$ for the integral of $y$ against the kernel $K$ on $[0,1]$ at $x$.

\section{Equations of the second kind}

\begin{lemma}\label{volterra}
$e$ is a solution of the Volterra integral equation with kernel $K$ and right-hand side $o$ from $0$.
\end{lemma}

\begin{lemma}\label{fredholm}
$c$ is a solution of the Fredholm integral equation with kernel $K$, right-hand side $o$ and parameter $\frac12$ on $[0,1]$, and $z$ is a solution of the homogeneous Fredholm integral equation with kernel $K$ and parameter $1$ on $[0,1]$.
\end{lemma}

\begin{lemma}\label{abel}
$o$ is a solution of Abel's integral equation with right-hand side $r$ from $0$.
\end{lemma}

\section{Kernels}

\begin{lemma}\label{kernels}
$K$ is a symmetric kernel, and $T(o,5)=1$.
\end{lemma}

\begin{lemma}\label{iterated}
Let $x,t\in\mathbb{R}$. Then $K_2(x,t)=1$.
\end{lemma}

\end{document}
